% Appendix: Prompt templates.
% Requires: \usepackage[most]{tcolorbox}
% All strings below are verbatim from the run configs / prompts.py.
% NB: the skill cards contain Trigger / Do / Avoid / Risk but NO explicit
% "Check" step, so the schema is stated as Trigger / Do / Avoid / Risk.

\subsection{Prompt templates}
\label{app:prompts}

Every condition shares one fixed system header and is assembled from two blocks,
a \emph{problem block} and a \emph{skill-card block}, concatenated in a
condition-specific order. We evaluate two headers, shown side by side in
Figure~\ref{fig:promptboxes}. The \textbf{TRS} header is the one used for every
result in the main analysis. The \textbf{directive} header is used only for the
prompt-sensitivity robustness check (Table~\ref{tab:promptsens}); it adds an
explicit instruction to exploit the hints and to keep the answer short.

\begin{figure*}[t]
\centering
\begin{minipage}[t]{0.485\textwidth}
\begin{tcolorbox}[colback=green!4, colframe=green!45!black,
    title=\textbf{TRS header (used for all analysis)}, fonttitle=\small,
    boxrule=0.6pt, left=3pt, right=3pt, top=2pt, bottom=2pt]
{\ttfamily\scriptsize
You are a helpful and harmless assistant.\\
You may be given an optional Solving Hints section. Use it only if it is relevant to the problem; otherwise, ignore it completely.\\
End your response with exactly one line in this format:\\
Final answer: \textless answer\textgreater
}
\end{tcolorbox}
\end{minipage}\hfill
\begin{minipage}[t]{0.485\textwidth}
\begin{tcolorbox}[colback=black!3, colframe=black!45,
    title=\textbf{Directive header (extended / robustness only)}, fonttitle=\small,
    boxrule=0.6pt, left=3pt, right=3pt, top=2pt, bottom=2pt]
{\ttfamily\scriptsize
You are a helpful and harmless assistant.\\
You may be given a Solving Hints section. Use it only if relevant; otherwise ignore it completely.\\
\textbf{If you find these solving hints useful, please try to reduce the number of tokens used while staying correct.}\\
End your response with exactly one line in this format:\\
Final answer: \textless answer\textgreater
}
\end{tcolorbox}
\end{minipage}
\caption{System headers. The two prompts differ only in the hint-usage
instruction: the directive header (right) adds the bold sentence and drops the
words ``optional'' and ``to the problem''. Everything else in the pipeline is
identical. All results outside Table~\ref{tab:promptsens} use the TRS header.}
\label{fig:promptboxes}
\end{figure*}

\paragraph{Blocks and separator.}
The two blocks are fixed across conditions:
%
\begin{center}
\ttfamily\small
\begin{tabular}{@{}ll@{}}
problem block: & Problem:\\
               & \{question\}\\[4pt]
skill block:   & [Solving Hints]\\
               & \{card\}\\
               & [/Solving Hints]\\
\end{tabular}
\end{center}
%
Consecutive blocks are joined by a blank line. When a segment is repeated, the
copies are separated by the token \texttt{Let me repeat it:} on its own line
(rendered with a blank line above and below). The header is prepended once, at
the top, in all conditions.

\paragraph{Conditions.}
Writing \texttt{q} for the problem block, \texttt{i} for the skill-card block,
and \texttt{|} for the separator token, the ten conditions are:
%
\begin{center}
\small
\begin{tabular}{@{}cll@{}}
\toprule
$k$ & sequence & description \\
\midrule
1  & \texttt{i\ q}          & insight-first, single \\
2  & \texttt{q\ i}          & problem-first, single \\
3  & \texttt{i\ q\ |\ i\ q} & insight-first, full pair repeated \\
4  & \texttt{q\ i\ |\ q\ i} & problem-first, full pair repeated \\
5  & \texttt{q}             & problem only (no card) \\
6  & \texttt{q\ |\ q}       & problem only, repeated \\
7  & \texttt{q\ i\ |\ q}    & problem-first $+$ problem-only bridge \\
8  & \texttt{q\ |\ q\ i}    & problem-only bridge $+$ problem-first \\
9  & \texttt{i\ q\ |\ i}    & insight-first $+$ card-only bridge \\
10 & \texttt{i\ |\ i\ q}    & card-only bridge $+$ insight-first \\
\bottomrule
\end{tabular}
\end{center}

\paragraph{Skill-card format.}
Each \texttt{\{card\}} is a retrieved DeepMath-103K OSS skill entry: a numbered
list of four to six heuristics. Every item follows an implicit
\emph{Trigger\,/\,Do\,/\,Avoid\,/\,Risk} schema, as in this verbatim example:
%
\begin{quote}\small
\emph{1. When encountering a difference of radicals with high indices as a
variable approaches infinity} [\textsc{Trigger}], \emph{adopt a Binomial
Expansion or Taylor Series approach} [\textsc{Do}] \emph{rather than
rationalization} [\textsc{Avoid}]. Rationalizing $n$-th roots becomes
algebraically prohibitive, whereas expanding $(x+a)^k$ quickly reveals the
leading order of growth. \emph{Be cautious of cases where the leading terms
cancel} [\textsc{Risk}].
\end{quote}
%
The cards are used exactly as retrieved; we do not reformat them into explicit
field labels.
